% ─────────────────────────────────────────────────────────────────────────────
% SHARED coursework preamble — every course inputs THIS.
%
% PREAMBLE ONLY. No \begin{document}, no course-specific notation.
%
% Course-specific macros live at COURSE level, in
%   content/<course>/reference_docs/<course>_macros.tex
% which is inputted AFTER this file. ONE per course, shared by hw/qz/exam --
% macros under a single kind are unreachable from a sibling kind.
%
% That split is the point: the page style, the answer box, and the problem
% header are identical everywhere, so they are defined once; \conv and \ustep
% mean something in DSP and nothing in a networks course, so they are not.
%
% A per-problem file therefore starts:
%   % ─────────────────────────────────────────────────────────────────────────────
% SHARED coursework preamble — every course inputs THIS.
%
% PREAMBLE ONLY. No \begin{document}, no course-specific notation.
%
% Course-specific macros live at COURSE level, in
%   content/<course>/reference_docs/<course>_macros.tex
% which is inputted AFTER this file. ONE per course, shared by hw/qz/exam --
% macros under a single kind are unreachable from a sibling kind.
%
% That split is the point: the page style, the answer box, and the problem
% header are identical everywhere, so they are defined once; \conv and \ustep
% mean something in DSP and nothing in a networks course, so they are not.
%
% A per-problem file therefore starts:
%   % ─────────────────────────────────────────────────────────────────────────────
% SHARED coursework preamble — every course inputs THIS.
%
% PREAMBLE ONLY. No \begin{document}, no course-specific notation.
%
% Course-specific macros live at COURSE level, in
%   content/<course>/reference_docs/<course>_macros.tex
% which is inputted AFTER this file. ONE per course, shared by hw/qz/exam --
% macros under a single kind are unreachable from a sibling kind.
%
% That split is the point: the page style, the answer box, and the problem
% header are identical everywhere, so they are defined once; \conv and \ustep
% mean something in DSP and nothing in a networks course, so they are not.
%
% A per-problem file therefore starts:
%   \input{../../../../docs/latex/coursework_preamble.tex}
%   \input{../../reference_docs/cesc470_macros.tex}
%   \begin{document}
%
% Build-check one file, or a whole assignment:
%   docs/latex/build_tex.sh content/cesc_470/hw/hw01/p01_five_components.tex
%   docs/latex/build_tex.sh content/cesc_470/hw/hw01
%
% Doctrine: docs/directives/coursework-solutions.md
% ─────────────────────────────────────────────────────────────────────────────

\documentclass[11pt]{article}

\usepackage[margin=1in]{geometry}
\usepackage{amsmath, amssymb, mathtools}
\usepackage{graphicx}
\usepackage{float}
\usepackage{booktabs}
\usepackage{longtable}
\usepackage{enumitem}
\usepackage{siunitx}
\usepackage[dvipsnames]{xcolor}
\usepackage{tcolorbox}
\usepackage{fancyhdr}
\usepackage{caption}
\usepackage{tikz}
\usepackage{pgfplots}
\pgfplotsset{compat=1.18}
\usepackage[colorlinks=true, allcolors=NavyBlue]{hyperref}

\captionsetup{font=small, labelfont=bf}
\setlist[enumerate]{itemsep=2pt, topsep=4pt}
\setlist[itemize]{itemsep=2pt, topsep=4pt}

% --- final-answer box -------------------------------------------------------
% Every problem file ends with its result in one of these. The condensed
% solutions document is assembled by lifting exactly these.
\newtcolorbox{answerbox}[1][]{
  colback=Green!6, colframe=Green!55!black,
  boxrule=0.6pt, arc=2pt, left=6pt, right=6pt, top=4pt, bottom=4pt,
  #1
}

% --- citation to course material --------------------------------------------
% \cite{Module 01, PDF p55}  ->  a small italic parenthetical.
% Solutions cite the SLIDE/PAGE a formula came from so a grader (or a future
% reader) can check the claim against what was actually taught. See the
% directive: every non-obvious step carries one.
\newcommand{\srcref}[1]{{\small\itshape(#1)}}

% --- a step that came from OUTSIDE the course materials ---------------------
% Used only where the course is genuinely silent and the problem asks for
% outside research. Marked visually so it is never mistaken for lecture content.
\newcommand{\extref}[1]{{\small\itshape\color{Sepia}[external: #1]}}

% --- callout for the trap in a problem --------------------------------------
\newtcolorbox{trap}[1][]{
  colback=BurntOrange!7, colframe=BurntOrange!70!black,
  boxrule=0.6pt, arc=2pt, left=6pt, right=6pt, top=4pt, bottom=4pt,
  fonttitle=\bfseries, title=Watch out, #1
}

% --- per-problem header -----------------------------------------------------
% \problemheader{8}{15 pts}{Target clock rate}
% The optional 4-argument form carries a learning outcome where the course
% states one (CESC 410 does; CESC 470's HW1 does not).
\newcommand{\problemheader}[3]{%
  \begin{center}\large\bfseries
    Problem #1 \quad\normalsize\mdseries(#2)\\[2pt]
    \large\bfseries #3
  \end{center}
  \vspace{2pt}\hrule\vspace{8pt}
}
\newcommand{\problemheaderlo}[4]{%
  \begin{center}\large\bfseries
    Problem #1 \quad\normalsize\mdseries(#2, #3)\\[2pt]
    \large\bfseries #4
  \end{center}
  \vspace{2pt}\hrule\vspace{8pt}
}

% Course name in the running head. Defined BEFORE \lhead references it, and
% \renewcommand'd by each course's macros file (inputted after this one).
\providecommand{\coursename}{}

\pagestyle{fancy}
\fancyhf{}
\lhead{\small\coursename}
\rhead{\small\thepage}
\renewcommand{\headrulewidth}{0.3pt}

%   % CESC 470 Computer Architecture — course-specific macros ONLY.
%
% Inputted AFTER docs/latex/coursework_preamble.tex, which carries everything
% shared (answerbox, problemheader, srcref, page style). Put nothing general
% here: if another course would want it, it belongs in the shared preamble.

\renewcommand{\coursename}{CESC 470}

% --- performance-equation notation ------------------------------------------
% The course writes these in words on the slides; these render them compactly
% and, more importantly, IDENTICALLY across every problem file.
\newcommand{\CPUtime}{T_{\mathrm{CPU}}}
\newcommand{\IC}{\mathrm{IC}}                    % instruction count
\newcommand{\CPI}{\mathrm{CPI}}
\newcommand{\cycletime}{t_{\mathrm{cycle}}}
\newcommand{\clockrate}{f_{\mathrm{clk}}}
\newcommand{\cycles}{N_{\mathrm{cycles}}}
\newcommand{\speedup}{S}

% Amdahl's law, in the form the slides state it (Module 01, PDF p86, slide 49).
\newcommand{\amdahl}[3]{#1 + \dfrac{#2}{#3}}

%   \begin{document}
%
% Build-check one file, or a whole assignment:
%   docs/latex/build_tex.sh content/cesc_470/hw/hw01/p01_five_components.tex
%   docs/latex/build_tex.sh content/cesc_470/hw/hw01
%
% Doctrine: docs/directives/coursework-solutions.md
% ─────────────────────────────────────────────────────────────────────────────

\documentclass[11pt]{article}

\usepackage[margin=1in]{geometry}
\usepackage{amsmath, amssymb, mathtools}
\usepackage{graphicx}
\usepackage{float}
\usepackage{booktabs}
\usepackage{longtable}
\usepackage{enumitem}
\usepackage{siunitx}
\usepackage[dvipsnames]{xcolor}
\usepackage{tcolorbox}
\usepackage{fancyhdr}
\usepackage{caption}
\usepackage{tikz}
\usepackage{pgfplots}
\pgfplotsset{compat=1.18}
\usepackage[colorlinks=true, allcolors=NavyBlue]{hyperref}

\captionsetup{font=small, labelfont=bf}
\setlist[enumerate]{itemsep=2pt, topsep=4pt}
\setlist[itemize]{itemsep=2pt, topsep=4pt}

% --- final-answer box -------------------------------------------------------
% Every problem file ends with its result in one of these. The condensed
% solutions document is assembled by lifting exactly these.
\newtcolorbox{answerbox}[1][]{
  colback=Green!6, colframe=Green!55!black,
  boxrule=0.6pt, arc=2pt, left=6pt, right=6pt, top=4pt, bottom=4pt,
  #1
}

% --- citation to course material --------------------------------------------
% \cite{Module 01, PDF p55}  ->  a small italic parenthetical.
% Solutions cite the SLIDE/PAGE a formula came from so a grader (or a future
% reader) can check the claim against what was actually taught. See the
% directive: every non-obvious step carries one.
\newcommand{\srcref}[1]{{\small\itshape(#1)}}

% --- a step that came from OUTSIDE the course materials ---------------------
% Used only where the course is genuinely silent and the problem asks for
% outside research. Marked visually so it is never mistaken for lecture content.
\newcommand{\extref}[1]{{\small\itshape\color{Sepia}[external: #1]}}

% --- callout for the trap in a problem --------------------------------------
\newtcolorbox{trap}[1][]{
  colback=BurntOrange!7, colframe=BurntOrange!70!black,
  boxrule=0.6pt, arc=2pt, left=6pt, right=6pt, top=4pt, bottom=4pt,
  fonttitle=\bfseries, title=Watch out, #1
}

% --- per-problem header -----------------------------------------------------
% \problemheader{8}{15 pts}{Target clock rate}
% The optional 4-argument form carries a learning outcome where the course
% states one (CESC 410 does; CESC 470's HW1 does not).
\newcommand{\problemheader}[3]{%
  \begin{center}\large\bfseries
    Problem #1 \quad\normalsize\mdseries(#2)\\[2pt]
    \large\bfseries #3
  \end{center}
  \vspace{2pt}\hrule\vspace{8pt}
}
\newcommand{\problemheaderlo}[4]{%
  \begin{center}\large\bfseries
    Problem #1 \quad\normalsize\mdseries(#2, #3)\\[2pt]
    \large\bfseries #4
  \end{center}
  \vspace{2pt}\hrule\vspace{8pt}
}

% Course name in the running head. Defined BEFORE \lhead references it, and
% \renewcommand'd by each course's macros file (inputted after this one).
\providecommand{\coursename}{}

\pagestyle{fancy}
\fancyhf{}
\lhead{\small\coursename}
\rhead{\small\thepage}
\renewcommand{\headrulewidth}{0.3pt}

%   % CESC 470 Computer Architecture — course-specific macros ONLY.
%
% Inputted AFTER docs/latex/coursework_preamble.tex, which carries everything
% shared (answerbox, problemheader, srcref, page style). Put nothing general
% here: if another course would want it, it belongs in the shared preamble.

\renewcommand{\coursename}{CESC 470}

% --- performance-equation notation ------------------------------------------
% The course writes these in words on the slides; these render them compactly
% and, more importantly, IDENTICALLY across every problem file.
\newcommand{\CPUtime}{T_{\mathrm{CPU}}}
\newcommand{\IC}{\mathrm{IC}}                    % instruction count
\newcommand{\CPI}{\mathrm{CPI}}
\newcommand{\cycletime}{t_{\mathrm{cycle}}}
\newcommand{\clockrate}{f_{\mathrm{clk}}}
\newcommand{\cycles}{N_{\mathrm{cycles}}}
\newcommand{\speedup}{S}

% Amdahl's law, in the form the slides state it (Module 01, PDF p86, slide 49).
\newcommand{\amdahl}[3]{#1 + \dfrac{#2}{#3}}

%   \begin{document}
%
% Build-check one file, or a whole assignment:
%   docs/latex/build_tex.sh content/cesc_470/hw/hw01/p01_five_components.tex
%   docs/latex/build_tex.sh content/cesc_470/hw/hw01
%
% Doctrine: docs/directives/coursework-solutions.md
% ─────────────────────────────────────────────────────────────────────────────

\documentclass[11pt]{article}

\usepackage[margin=1in]{geometry}
\usepackage{amsmath, amssymb, mathtools}
\usepackage{graphicx}
\usepackage{float}
\usepackage{booktabs}
\usepackage{longtable}
\usepackage{enumitem}
\usepackage{siunitx}
\usepackage[dvipsnames]{xcolor}
\usepackage{tcolorbox}
\usepackage{fancyhdr}
\usepackage{caption}
\usepackage{tikz}
\usepackage{pgfplots}
\pgfplotsset{compat=1.18}
\usepackage[colorlinks=true, allcolors=NavyBlue]{hyperref}

\captionsetup{font=small, labelfont=bf}
\setlist[enumerate]{itemsep=2pt, topsep=4pt}
\setlist[itemize]{itemsep=2pt, topsep=4pt}

% --- final-answer box -------------------------------------------------------
% Every problem file ends with its result in one of these. The condensed
% solutions document is assembled by lifting exactly these.
\newtcolorbox{answerbox}[1][]{
  colback=Green!6, colframe=Green!55!black,
  boxrule=0.6pt, arc=2pt, left=6pt, right=6pt, top=4pt, bottom=4pt,
  #1
}

% --- citation to course material --------------------------------------------
% \cite{Module 01, PDF p55}  ->  a small italic parenthetical.
% Solutions cite the SLIDE/PAGE a formula came from so a grader (or a future
% reader) can check the claim against what was actually taught. See the
% directive: every non-obvious step carries one.
\newcommand{\srcref}[1]{{\small\itshape(#1)}}

% --- a step that came from OUTSIDE the course materials ---------------------
% Used only where the course is genuinely silent and the problem asks for
% outside research. Marked visually so it is never mistaken for lecture content.
\newcommand{\extref}[1]{{\small\itshape\color{Sepia}[external: #1]}}

% --- callout for the trap in a problem --------------------------------------
\newtcolorbox{trap}[1][]{
  colback=BurntOrange!7, colframe=BurntOrange!70!black,
  boxrule=0.6pt, arc=2pt, left=6pt, right=6pt, top=4pt, bottom=4pt,
  fonttitle=\bfseries, title=Watch out, #1
}

% --- per-problem header -----------------------------------------------------
% \problemheader{8}{15 pts}{Target clock rate}
% The optional 4-argument form carries a learning outcome where the course
% states one (CESC 410 does; CESC 470's HW1 does not).
\newcommand{\problemheader}[3]{%
  \begin{center}\large\bfseries
    Problem #1 \quad\normalsize\mdseries(#2)\\[2pt]
    \large\bfseries #3
  \end{center}
  \vspace{2pt}\hrule\vspace{8pt}
}
\newcommand{\problemheaderlo}[4]{%
  \begin{center}\large\bfseries
    Problem #1 \quad\normalsize\mdseries(#2, #3)\\[2pt]
    \large\bfseries #4
  \end{center}
  \vspace{2pt}\hrule\vspace{8pt}
}

% Course name in the running head. Defined BEFORE \lhead references it, and
% \renewcommand'd by each course's macros file (inputted after this one).
\providecommand{\coursename}{}

\pagestyle{fancy}
\fancyhf{}
\lhead{\small\coursename}
\rhead{\small\thepage}
\renewcommand{\headrulewidth}{0.3pt}

% CESC 470 Computer Architecture — course-specific macros ONLY.
%
% Inputted AFTER docs/latex/coursework_preamble.tex, which carries everything
% shared (answerbox, problemheader, srcref, page style). Put nothing general
% here: if another course would want it, it belongs in the shared preamble.

\renewcommand{\coursename}{CESC 470}

% --- performance-equation notation ------------------------------------------
% The course writes these in words on the slides; these render them compactly
% and, more importantly, IDENTICALLY across every problem file.
\newcommand{\CPUtime}{T_{\mathrm{CPU}}}
\newcommand{\IC}{\mathrm{IC}}                    % instruction count
\newcommand{\CPI}{\mathrm{CPI}}
\newcommand{\cycletime}{t_{\mathrm{cycle}}}
\newcommand{\clockrate}{f_{\mathrm{clk}}}
\newcommand{\cycles}{N_{\mathrm{cycles}}}
\newcommand{\speedup}{S}

% Amdahl's law, in the form the slides state it (Module 01, PDF p86, slide 49).
\newcommand{\amdahl}[3]{#1 + \dfrac{#2}{#3}}


\begin{document}

\problemheader{7}{25 pts}{ARM vs x86 --- advantages and disadvantages}

\subsection*{Problem statement}

ARM and x86 are two widely used ISAs. Research and compare one advantage and one
disadvantage of each in terms of performance, power consumption, or
compatibility.

\subsection*{Approach}

This is the only problem on the assignment that asks for outside research, and
it carries a quarter of the marks. The course names both ISAs
\srcref{Module 01, PDF p25, slide 12} but does not compare them, so the
substance must come from outside — marked as external throughout, per the
citation policy, so lecture content and external claims stay distinguishable.

Structure it as the question asks: for \textbf{each} ISA, one advantage and one
disadvantage, each tied to one of the three named axes (performance, power,
compatibility). Four claims, argued rather than asserted.

The organising idea is the historical design split, which explains every point
below:

\begin{center}
\begin{tabular}{lll}
  \toprule
  & \textbf{ARM} & \textbf{x86} \\
  \midrule
  Design lineage    & RISC — reduced instruction set & CISC — complex instruction set \\
  Instruction width & fixed (32-bit, or 16/32 in Thumb) & \textbf{variable}, 1--15 bytes \\
  Original target   & embedded, battery-powered & desktop performance \\
  Business model    & licensed IP, many implementers & implemented by Intel and AMD \\
  \bottomrule
\end{tabular}
\end{center}

\extref{ARM Ltd., \emph{ARM Architecture Reference Manual}; Intel,
\emph{Intel 64 and IA-32 Architectures Software Developer's Manual}, Vol.~2}

\subsection*{Work}

\subsubsection*{ARM}

\paragraph{Advantage --- power efficiency (power consumption).}
{\sloppy
ARM's fixed-width, regularly-encoded instructions make the \emph{decoder} small
and simple: the processor knows where each instruction begins without first
decoding the previous one, so it can fetch and decode several in parallel with
little logic. Fewer transistors switching per instruction means less energy per
instruction.\par}

This is why ARM dominates battery-powered devices, and why the design point
matters more than any single benchmark: performance-per-watt, not peak
performance, is the metric a phone or a laptop is judged on. It is also the
reason ARM parts scale down to microcontrollers drawing milliwatts, a range x86
does not serve.

\extref{Furber, \emph{ARM System-on-Chip Architecture}, 2nd ed., Addison-Wesley}

\paragraph{Disadvantage --- software compatibility and ecosystem fragmentation.}
Decades of desktop, server, and professional software was compiled for x86.
Running it on ARM requires recompilation (needs source) or binary translation
(costs performance and can break edge cases). Apple's Rosetta 2 and Microsoft's
x86-on-ARM emulation both demonstrate the transition is possible \emph{and} that
it is a real cost, not a free one.

A second compatibility wrinkle is specific to ARM's licensing model: because
many vendors implement the architecture, ARM systems have historically lacked a
single standardised platform/boot interface of the kind PCs have, so an OS image
often has to be built per device. \extref{Arm SystemReady / SBSA programme
documentation} Note this is a \emph{platform} issue rather than an ISA one
strictly speaking, but it is a genuine compatibility cost of the ecosystem.

\subsubsection*{x86}

\paragraph{Advantage --- backward compatibility (compatibility).}
x86 has maintained binary compatibility since the 8086 in 1978. Software
compiled decades ago still runs on current processors. For enterprises with
large bodies of legacy or closed-source software this is worth more than any
efficiency gain, because the alternative is rewriting or abandoning working
systems.

This is exactly the trade-off the course names as the value of the ISA
abstraction \srcref{PDF p25, slide 12}: ``advantage: different implementations
of the same architecture'' — x86 is the extreme case, with the same architecture
implemented across nearly fifty years of organizations.

\paragraph{Disadvantage --- decoding complexity, and the power cost of it (power / performance).}
x86 instructions are \textbf{variable length}, from 1 to 15 bytes. A decoder
cannot know where instruction $n+1$ starts until it has decoded instruction $n$,
which makes parallel decode genuinely hard and forces significant extra logic
(length pre-decode, complex multi-way decoders, micro-op caches). Modern x86
cores internally translate these instructions into RISC-like micro-operations —
an entire extra pipeline stage that exists purely to service the legacy
encoding.

That decode machinery costs transistors and power on every instruction, which is
the same conclusion the course draws in general terms when it explains why the
ISA abstraction can hurt: it ``sometimes prevents using new innovations''
\srcref{PDF p25, slide 12}. x86 cannot drop the legacy encoding without breaking
the compatibility that is its main advantage.

\extref{Intel SDM Vol.~2 (instruction encoding); Hennessy \& Patterson,
\emph{Computer Architecture: A Quantitative Approach}, App.~A on ISA design}

\begin{trap}
\textbf{Do not overstate the ISA's role in performance.} It is tempting to
conclude ``RISC is faster'' or ``CISC is slower''. Measured performance is
dominated by \emph{organization} — pipeline depth, cache hierarchy, issue width,
process node — not by the ISA \srcref{PDF p31, slide 16}. High-performance x86
and high-performance ARM parts are competitive with each other; the durable ISA
differences are in \emph{decode complexity} and \emph{energy per instruction},
which is why the answer above is framed that way rather than as raw speed.
\end{trap}

\subsection*{Check}

Each of the four claims is tied to one of the three axes the question names, and
each traces to the same root cause, which is the sign the comparison is coherent
rather than a list of assorted facts:

\begin{center}
\begin{tabular}{llll}
  \toprule
  ISA & & Axis & Root cause \\
  \midrule
  ARM & advantage    & power         & fixed-width encoding $\to$ simple decoder \\
  ARM & disadvantage & compatibility & later arrival to the desktop/server software base \\
  x86 & advantage    & compatibility & unbroken binary compatibility since 1978 \\
  x86 & disadvantage & power / perf  & variable-length encoding $\to$ complex decoder \\
  \bottomrule
\end{tabular}
\end{center}

Note the symmetry: each ISA's advantage is the direct cause of the other's
disadvantage. ARM's simple encoding buys efficiency but no legacy base; x86's
legacy base is preserved by an encoding that costs efficiency. That is a
trade-off, not a ranking — which is what the question is really testing.

\subsection*{Answer}

\begin{answerbox}
\textbf{ARM}
\begin{itemize}[nosep, leftmargin=1.4em]
  \item \textbf{Advantage --- power consumption.} Fixed-width, regular RISC
    encoding keeps the instruction decoder small and simple, so less energy is
    spent per instruction. This is why ARM dominates battery-powered devices,
    where performance-per-watt is the metric that matters.
  \item \textbf{Disadvantage --- compatibility.} The large body of existing
    desktop/server software was compiled for x86 and must be recompiled or
    emulated (Rosetta 2, Windows-on-ARM), which costs performance and can break
    edge cases.
\end{itemize}

\medskip
\textbf{x86}
\begin{itemize}[nosep, leftmargin=1.4em]
  \item \textbf{Advantage --- compatibility.} Binary compatibility maintained
    since the 8086 (1978): decades-old software still runs on current
    processors, which is decisive for legacy and closed-source codebases.
  \item \textbf{Disadvantage --- decode complexity and its power cost.}
    Variable-length instructions (1--15 bytes) mean the start of instruction
    $n+1$ is unknown until instruction $n$ is decoded. Parallel decode therefore
    needs substantial extra logic, and cores must translate to internal
    micro-operations — transistors and power spent servicing a legacy encoding
    that cannot be dropped without losing the compatibility above.
\end{itemize}

\medskip
\textbf{Caveat:} raw performance is governed mainly by \emph{organization}
(pipeline, caches, process node), not by the ISA \srcref{PDF p31}. The durable
ISA-level differences are decode complexity and energy per instruction.
\end{answerbox}

\end{document}
